\section{Contexto largo, memoria y aprendizaje continuo: los límites de la capacidad se están redibujando}

\subsection{Avances en Context Length y la «ansiedad por la compresión»}

La ventana de contexto sigue creciendo, pero la entrevista subraya que un contexto largo no es gratuito. Una vez que la ventana se agranda, el nuevo problema es cómo recuperar la información clave, cómo comprimir sin perder el significado y cómo mantener la coherencia a lo largo de cadenas largas. En muchas herramientas, la «compresión automática» reduce detalles de alto valor a resúmenes toscos, lo que distorsiona el razonamiento posterior.

\begin{figure}[H]
\centering
\includegraphics[width=0.88\textwidth]{frame-03.jpg}
\caption{Discusión de la entrevista sobre el contexto largo y las estrategias de compresión\protect\footnotemark}
\end{figure}
\footnotetext{Intervalo de tiempo de la imagen del video: 02:44:10--02:44:20.}

\begin{importantbox}{Tres claves de ingeniería para los sistemas de contexto largo}
\begin{itemize}[leftmargin=1.5em]
    \item \textbf{Estrategia de recuperación de información}: no se trata de «dárselo todo al modelo», sino de filtrar antes de leer;
    \item \textbf{Estrategia de compresión}: conservar los hechos accionables y descartar la retórica de bajo valor;
    \item \textbf{Estrategia de restauración}: cuando el resumen se distorsiona, poder retroceder al fragmento original y corregir el sesgo.
\end{itemize}
\end{importantbox}

\subsection{Continual Learning: ¿actualización de pesos o inyección de contexto?}

La entrevista divide el aprendizaje continuo en dos caminos: actualizar los pesos (el aprendizaje propiamente dicho) e inyectar memoria (un complemento en tiempo de ejecución). El primero tiene un costo y un riesgo altos, pero es más estable a largo plazo; el segundo se implementa rápido y cuesta poco, pero está sujeto a las limitaciones de la ventana de contexto y de la calidad de la recuperación. Los sistemas reales suelen adoptar una ruta híbrida.

\begin{knowledgebox}{LoRA y el equilibrio real de la actualización en línea}
Métodos de adaptación de bajo rango como LoRA son comunes en la personalización empresarial porque ofrecen un equilibrio controlable entre la «velocidad de aprendizaje» y el «riesgo de olvido». Si se busca una actualización en línea a escala de minutos, se necesita una capacidad sólida de monitoreo y reversión; de lo contrario, es muy fácil que la retroalimentación con ruido termine escrita en el comportamiento del modelo.
\end{knowledgebox}

\begin{warningbox}{«Recordar al usuario» no equivale a «entender al usuario»}
Acumular el historial del usuario en el contexto puede mejorar la personalización, pero también puede amplificar sesgos antiguos, introducir riesgos de filtración de privacidad y hacer que el modelo dependa en exceso de plantillas viejas en tareas nuevas. Los sistemas de personalización deben contar con límites claros de memoria y mecanismos de eliminación.
\end{warningbox}

\subsection{Resumen de la sección}

El contexto largo y el aprendizaje continuo son variables clave del límite superior de la capacidad, pero el verdadero problema está en el control de ingeniería: la recuperación de información, la compresión, la reversión, la privacidad y la evaluación deben diseñarse a la vez.

